\documentclass[11pt]{article}

\usepackage[margin=1in]{geometry}
\usepackage{amsmath,amssymb}
\usepackage{booktabs}
\usepackage{array}
\usepackage{textcomp}
\usepackage[hidelinks]{hyperref}
\usepackage{listings}
\lstdefinestyle{cfg}{basicstyle=\ttfamily\footnotesize, breaklines=true,
  columns=fullflexible, keepspaces=true, showstringspaces=false,
  frame=single, framesep=5pt, xleftmargin=4pt, aboveskip=6pt, belowskip=4pt}

\newcommand{\code}[1]{\texttt{#1}}
\newcommand{\eps}{\varepsilon}

\title{\textbf{Wirewalker Data Processing: Methods and Implementation}}
\author{PyWirewalker}
\date{\today}

\begin{document}
\maketitle

\begin{abstract}
\noindent
PyWirewalker provides processing pipelines for the instruments carried on a wave-powered
Wirewalker profiling mooring. This report documents the processing architecture, the RBR
Concerto CTD chain, and the Nortek Signature1000 ADCP chain, with particular emphasis on the
improvements introduced when the ADCP toolbox was ported from its original MATLAB
implementation to Python. The code, configuration templates, and this report are openly
available at \url{https://github.com/missioncreep11/PyWirewalker}.
\end{abstract}

\section{Processing architecture (L0--L3)}

The pipeline follows a staged \textbf{level (L0--L3)} convention in which each level is a
self-describing NetCDF product derived deterministically from the level below it. Raw
instrument files are never modified; every derived quantity is reproducible from the raw
data and a single version-controlled configuration file.

\begin{table}[h]
\centering
\begin{tabular}{@{}p{0.07\textwidth} p{0.42\textwidth} p{0.40\textwidth}@{}}
\toprule
\textbf{Level} & \textbf{Content} & \textbf{Character} \\
\midrule
\textbf{L0} & Raw instrument file (RBR \code{.rsk}, Nortek \code{.ad2cp}) & Vendor binary; read-only, never edited \\
\textbf{L1} & Full-resolution converted time series & Physical units, per-sample cast/profile flags \\
\textbf{L2} & Gridded per-cast product, dimensioned \code{(depth, cast)} & QC applied, derived variables, one profile per column \\
\textbf{L3} & Regular \code{(depth, time)} grid & Uniform time base for spectral/temporal analysis \\
\bottomrule
\end{tabular}
\caption{Processing levels.}
\end{table}

Two design principles are common to both instruments:

\begin{enumerate}
\item \textbf{Strict downward dependence.} $\mathrm{L}(n)$ reads only $\mathrm{L}(n-1)$ ---
never the raw file --- so a given level can inherit only the variables present in the level
beneath it. This makes the provenance of every field explicit and prevents silent
reprocessing from source.

\item \textbf{Configuration, not code, carries deployment specifics.} All deployment- and
machine-specific settings (paths, metadata, calibration constants, grid parameters) live in
a JSON configuration file --- \code{config\_ctd.json} for the CTD, \code{config\_adcp.json}
for the ADCP. No paths or constants are hard-coded; the resolved values are written into the
output NetCDF attributes for provenance. A single driver processes any deployment by pointing
it at the appropriate config (\code{--config}), with individual settings overridable on the
command line. Consequently, every deployment-specific quantity quoted in this report ---
sample rate, grid spacings, correction parameters, QC thresholds, and metadata --- is a
configuration default for the reference deployment, and the governing key is named where each
value first appears.
\end{enumerate}

The \code{(depth, cast)} convention for L2 is shared by the CTD and ADCP: each column is one
profile (upcast), each row a depth bin, and a \code{time} coordinate carries the cast
mid-time. This common structure lets CTD, velocity, and turbulence products be co-registered
by cast and depth without interpolation. The two instrument chains are otherwise independent
--- different raw formats, different drivers, different physics --- and are described in turn
below.

\section{CTD processing --- RBR Concerto}

The CTD chain (\code{ww\_rbr/}, driver \code{process\_wirewalker\_rbr.py}) converts a raw RBR
Concerto \code{.rsk} file into an L1--L3 archive. The reference deployment is mooring
\textbf{NOPP-Aleutians} (RBR Concerto\textsuperscript{3} S/N 213752, 2~Hz continuous,
$\approx 5.3\times10^{7}$ scans over 2025-07 to 2026-05).

\subsection{Cast detection}

Profiles are not re-detected; the chain reuses the instrument-generated segmentation stored
in the Ruskin \code{region}\,/\,\code{regionCast} tables. Each profile is split into a
\emph{down cast} (the slow, ratcheting descent as the vehicle sinks) and an \emph{up cast}
(the fast, buoyant ascent). Only \textbf{upcasts} are retained for gridded products: the CTD
is mounted on top of the vehicle and sits in its turbulent wake during descent, so downcasts
are systematically contaminated. Up/down agreement is therefore never used as a quality
metric.

\subsection{L1 --- full-resolution conversion}

\code{build\_L1} reads the raw \code{.rsk} (a SQLite database) and writes the full-rate time
series in physical units at whatever rate the instrument was configured to sample. The
sample rate is not assumed anywhere in the chain --- it is read from the \code{.rsk} and
carried through as the \code{sampling\_hz} configuration item (2~Hz for the reference
deployment above, but arbitrary in general; other deployments in this project run at 6 and
8~Hz). The measured channels are \textbf{discovered from the
instrument metadata tables} (\code{channels}, \code{instrumentChannels}) rather than
hard-coded, so auxiliary sensors (optical backscatter, chlorophyll fluorescence, dissolved
oxygen, PAR, a second thermistor, \dots) are carried through automatically under names
slugged from their long names. Three channels receive canonical names for the downstream
TEOS-10 step: conductivity, the \code{temp14} C--T-cell thermistor, and total pressure. Every
sample is tagged with \code{cast\_number}, \code{profile\_number}, and \code{cast\_direction}
($0=$ down, $1=$ up). Conductivity is stored raw at L1; no thermal-mass correction is applied
at this stage.

\subsection{L2 --- gridded upcasts, thermodynamics, and de-spiking}

\code{build\_L2} reads the \textbf{L1 NetCDF} (not the raw file), selects upcasts, and
bin-averages each onto a vertical grid whose bin size and extent are config items
(\code{l2\_dz\_m}, \code{zmin\_m}, \code{zmax\_m}; 0.5~m bins over 0--500~m for the reference
deployment). Two corrections are
applied per cast before gridding:

\begin{itemize}
\item \textbf{Conductivity-cell thermal-mass correction} (Lueck \& Picklo, 1990), addressing
the thermal inertia of the conductivity cell that otherwise produces salinity spikes at sharp
thermal gradients. The recursive correction uses the parameters in the config \code{thermal\_mass} block, which
default to the RBR \code{pyRSKtools} values $\alpha = 0.04$, $\beta = 0.1~\mathrm{s^{-1}}$
($\tau \approx 10$~s), $\gamma = 1.0$. Fitting
$\tau$ freely is ill-posed on these profiles (the estimate runs away and biases the profile),
so the fixed manufacturer parameters are used and recorded in the L2 attributes.

\item \textbf{Conductivity--temperature alignment.} The optimal C--T lag was measured at
$\approx 0$~s across the deployment (the Concerto's C and T channels are already aligned), so
no lag shift is applied.
\end{itemize}

Thermodynamic variables are then derived with the TEOS-10 Gibbs SeaWater toolbox
(\code{gsw}): sea pressure $p = p_{\mathrm{total}} - p_{\mathrm{atm}}$ (with the atmospheric
offset $p_{\mathrm{atm}}$ set in the config, \code{atmospheric\_pressure\_dbar}, defaulting to
the \code{.rsk} \code{ATMOSPHERE} value of 10.1325~dbar), depth from \code{gsw.z\_from\_p},
practical salinity \code{SP\_from\_C}, absolute salinity, conservative temperature, potential
density anomaly $\sigma_0$, and sound speed --- the latitude and longitude required by the gsw
conversions likewise taken from the config (\code{latitude}, \code{longitude}). The L2 product carries
these together with the bin-averaged auxiliary channels and a per-bin sample count
\code{n\_obs}, dimensioned \code{(depth, cast)}.

\subsection{L3 --- regular depth--time grid and stratification}

\code{build\_L3} reads L2 and interpolates onto a regular depth--time grid whose vertical and
temporal spacing are set in the configuration file (\code{l3\_dz\_m} and \code{l3\_dt};
\textbf{1~m $\times$ 30~min} for the reference deployment, chosen to match its cast cadence). The native upcast cadence ($\approx 31$~min) gives roughly one profile per
time bin (Nyquist $\approx 1$~cph); $\approx 88\%$ of bins contain a single upcast,
$\approx 4\%$ two or more, and $\approx 8\%$ are empty. Empty bins are left \textbf{NaN --- no
temporal interpolation} --- in the archival product to preserve sample independence; a
companion \code{\_interp} product fills only isolated single-bin gaps
($\le$ the config \code{l3\_interp\_max\_gap\_bins}) by linear interpolation in time, flagging filled bins with
\code{n\_casts == 0}.

The squared buoyancy frequency is computed on the L3 grid as
\begin{equation}
N^{2} = \frac{g}{\rho_0}\,\frac{\partial \sigma_0}{\partial z},
\end{equation}
with $z$ positive down (so $N^2 > 0$ is stable) and $g$ the config \code{gravity}, using the
gridded $\sigma_0$ field after a NaN-aware boxcar vertical smooth whose length is the config
\code{n2\_vertical\_smoothing\_m} (5~m for the reference deployment, selected from a
gradient-smoothing comparison that removes most differentiation noise while preserving the
pycnocline peak, and recorded in the product attributes).

\section{ADCP processing --- Nortek Signature1000}

The ADCP chain (\code{ww\_sig1000/}, driver \code{process\_ww\_sig1000.py}) processes the raw
Nortek \code{.ad2cp} file into two \code{(depth, cast)} L2 products: \textbf{velocity}
(motion-corrected ENU currents and vertical shear from the four slant beams) and
\textbf{turbulence} (spectral and structure-function dissipation $\eps$ from the fifth,
high-resolution beam). This chain is a Python port of a MATLAB toolbox (Zheng, Lucas,
Le~Boyer, and Northcott; upstream \code{modscripps/wirewalker}). The remainder of this section
describes the port and, throughout, highlights where it corrects or extends the original.

\subsection{Ingest and streaming architecture}

\textbf{Improvement --- direct raw ingest.} The original workflow required an intermediate
Nortek \code{.mat} export followed by bespoke \code{sort\_file}\,/\,\code{merge\_signature}
reassembly steps. The port reads the raw \code{.ad2cp} directly through MHKiT/DOLfYN
(\code{from mhkit import dolfyn}), which parses the native binary and exposes the slant-beam
burst, the fifth-beam high-resolution (HR) burst, the IMU/AHRS records, and the instrument
geometry (cell size, blanking distance, ambiguity velocity, beam angle, sample rate). The
export and reassembly stages are thereby eliminated, removing a lossy, manual step from the
pipeline.

\textbf{Improvement --- bounded-memory streaming.} Deployment files routinely exceed 10~GB
(the reference NOPP file is $\approx 38$~GB\,/\,$4\times10^{7}$ ensembles). The driver streams
the file in ensemble chunks (the chunk size is the config \code{chunk}), detecting casts on a
rolling pressure buffer and carrying any cast that spans a chunk boundary into the next read,
so memory use is bounded by the chunk size rather than the file size. Optional bounds set in
the config (\code{start\_time}/\code{end\_time}, or their ensemble equivalents) trim
deployment/recovery transit. Instrument geometry is read from the file itself and is therefore
never a configuration item.

\subsection{Velocity and shear (slant beams)}

Currents are formed by the standard beam $\to$ instrument (XYZ) $\to$ Earth (ENU)
transformation of the four slant beams, after a correlation-threshold mask (config
\code{corr\_min}), depth-alignment of the per-beam cells using the tilt-corrected cell
geometry, and box-averaging onto a uniform depth grid whose bin size and maximum depth are
config items (\code{boxsize\_m}, \code{z\_max\_m}; 1~m by default). The casts gridded are those
selected by direction and minimum pressure span (config \code{kind} and \code{min\_span\_dbar}).
Two motion-correction models are provided:

\begin{itemize}
\item \textbf{v1} --- a port of the original \code{WWcorr\_beam}: bandpass-integrated IMU
acceleration plus the pressure-rate ($dp/dt$) ascent, rotated by the AHRS attitude.
\item \textbf{v2 (improvement)} --- a buoyant-ascent model in which the vertical velocity is
taken from $dp/dt$, the horizontal platform motion is depth-gated, and the tilt is derived
from low-passed accelerometer data. This model is immune to AHRS attitude faults and adds a
\textbf{sail correction} for the along-wire travel of the vehicle on an inclined mooring line.
The attitude source is selectable (\code{ahrs}, \code{reconstructed}, or \code{auto}, which
reconstructs only on AHRS-faulted casts) and is recorded per cast. The model (\code{motion}),
whether any motion correction is applied (\code{motion\_correct}), and the attitude source
(\code{attitude}) are all configuration selections.
\end{itemize}

\textbf{Improvement --- motion-immune shear.} In addition to gridded velocity, the L2 product
carries a \textbf{beam-differenced vertical shear} (\code{shearE}, \code{shearN}), a port of
the toolbox's \code{beamshear}. Centred cell differences
$(v_{c+1}-v_{c-1})/(z_{c+1}-z_{c-1})$ are formed \textbf{along each raw beam before rotation}.
Any quantity common to a ping's cells --- platform translation, attitude-error leakage of the
ascent, the sail term --- cancels exactly in the difference, so the beam shear is immune to
the entire motion/attitude error family and requires no motion correction. The four beam
shears are then rotated to ENU (differentiation being linear) and box-averaged. Each velocity
and shear bin carries a Doppler-noise standard error (\code{\_sem}), enabling quantitative
signal-to-noise assessment (Section~\ref{sec:qc}).

\subsection{Turbulence --- spectral method (HR beam 5)}

Dissipation is estimated from the pulse-coherent HR fifth beam, which points along the
vertical and resolves fine-scale vertical velocity. The per-cast pipeline is:

\begin{enumerate}
\item Mask low-correlation samples (config \code{corr\_min}); convert velocity to phase,
unwrap, and remove the ensemble mean rotation (\emph{angular demean}), which suppresses
alias-correction spikes.
\item Compute each cell's vertical position $z = p - b_Z\,r$ from the tilt-corrected beam
geometry.
\item Estimate and remove the \emph{stagnation profile}: the median velocity profile over the
deep (high-pressure) portion of the cast, subtracted from every ping; the near-transducer
cutoff is the first cell where the deep median returns to within $4~\mathrm{mm\,s^{-1}}$ of its
own median, selected nearest 1~m range.
\item De-spike (single points with large $|dv/dz|$ on both sides) and detrend in range.
\item Form band-averaged vertical-wavenumber spectra $S(k)$ over depth bins (bin resolution
\code{dep\_res\_m} and maximum depth \code{max\_dep\_m} from the config) and fit the Kolmogorov
inertial-subrange model
\begin{equation}
S(k) = N + A\,k^{-5/3}, \qquad
\eps = \left(\frac{A}{C_K}\right)^{3/2},\quad C_K = 0.53,
\end{equation}
where $N$ is the (white) spectral noise floor and $A$ the inertial-subrange amplitude.
\end{enumerate}

\textbf{Improvement --- correct reference implementation.} The vendored script
\code{WWturb\_upward.m} had diverged from the method actually published by the instrument
authors and produced a systematic \textbf{$+0.28$~dex} high bias in $\eps$. The port instead
reproduces \code{ProcessSingleProfile.m} --- the final paper code from Northcott et al.\ (Dryad
\href{https://doi.org/10.5061/dryad.8sf7m0d44}{\code{doi:10.5061/dryad.8sf7m0d44}}). The key
corrections relative to the vendored script are: band-averaged spectra (rather than
interpolated), the data-dependent wavenumber-grid length, the deep-profile stagnation
subtraction (absent in the vendored code), the cutoff-nearest-1~m rule with a pressure-based
deep threshold, and removal of an extraneous averaging mask.

\textbf{Validation.} Against Northcott et al.'s published product
\code{NortekTurbulenceData.nc} for the full TLC 2023 deployment (2116 profiles $\times$ 67
depth bins at 3~m resolution), the ported spectral $\eps$ reproduces the published values with
a \textbf{median offset of $-0.00$~dex, RMS $\approx 0.08$~dex, and correlation of
$\log\eps$ of $0.994$}, and yields exactly the published profile count. The gold-standard
single profile (\code{RawProfile2\_20231005-1132.nc}) matches to $-0.00$~dex. This
bit-for-bit agreement anchors the port to the peer-reviewed result.

\subsection{Turbulence --- structure-function method (new)}

\textbf{Improvement --- a second, gap-robust estimator.} The port additionally implements the
second-order velocity structure-function method (a port of the SF branch of
\code{WWturb\_upward.m}) as a parallel estimate computed from the same preprocessed velocity.
The structure function is accumulated over both sides of each depth bin, counting only valid
sample pairs, and fit as
\begin{equation}
D(r) = N + A\,r^{2/3}, \qquad
\eps = \left(\frac{A}{C_{SF}}\right)^{3/2}.
\end{equation}
The absolute constant $C_{SF}$ is cross-calibrated to the validated spectral $\eps$ on the
high-scattering TLC deployment ($C_{SF} = 1.476$; the literature value is $\approx 2.0$), so
that the two estimators share one absolute scale. The spectral and SF estimates
(\code{epsilon}, \code{epsilon\_sf}) are stored side by side in the product.

The value of the second estimator is specific to \textbf{low-scattering environments}. The
spectral method zero-fills correlation-masked (gap) samples before the FFT; in high-return
water this is negligible ($\approx 8\%$ of cells masked at TLC), but in low-return water it
injects spurious broadband variance that raises the apparent noise floor. The structure
function, computed from velocity differences over valid pairs only, \textbf{skips gaps rather
than filling them}. On a low-scattering deployment (NOPP, beam-5 correlation $\approx 68\%$,
$\approx 30$--$45\%$ of cells masked), the two estimators agree only weakly, and the difference
is diagnostic --- see the quality framework below.

\subsection{Quality framework and diagnostics (new)}
\label{sec:qc}

The port introduces several quantitative quality measures that were not part of the original
toolbox.

\paragraph{$\eps$--noise-floor coupling.} Genuine turbulence should vary independently of the
instrument noise floor; where $\eps$ instead tracks $N$, it is being set by the instrument
rather than the ocean. The per-bin correlation $r(\log\eps, \log N)$ therefore serves as a
contamination diagnostic. Calibration on the validated (high-scattering) TLC deployment gives
$r \approx -0.3$ --- the mild \emph{negative} value expected from the intrinsic $N$--$A$
trade-off of the two-parameter fit --- with no diel structure. The low-scattering NOPP
deployment instead shows $r$ rising to $+0.6$ to $+0.9$ with depth, and a diel $\eps$ cycle in
the upper ocean \textbf{locked to the noise floor} ($r \approx 0.95$), identifying a biological
(diel vertical migration) contamination of the near-surface estimate and an increasingly
noise-limited deep estimate. The structure-function estimator \textbf{removes the gap-driven
component} of this coupling in mid-water (driving $r$ from $\approx +0.45$ to $\approx 0$),
extending the trustworthy depth range; the residual deep coupling reflects genuine biological
scattering, which no estimator can remove.

\paragraph{Transfer-function noise floor for shear.} A flat (white) noise floor derived from
the per-bin SEM is \emph{not} the correct floor for beam-differenced shear: the differencing
operator $(v_{c+1}-v_{c-1})$ and the box-average shape the noise spectrally. The correct
vertical-wavenumber noise floor is the white per-sample Doppler variance passed through the
processing transfer function,
\begin{equation}
|H(m)|^{2} \propto \sin^{2}\!\left(2\pi m\,c_z\right)\,
\operatorname{sinc}^{2}\!\left(m L\right),
\end{equation}
with $c_z$ the vertical cell projection and $L$ the box-average length, normalized so its
integral equals the SEM variance. Because differencing removes the mean, this floor rises as
$m^{2}$ at low wavenumber --- the Doppler noise is negligible at large vertical scales. On
NOPP this places the shear signal-to-noise above unity out to vertical scales of
$\approx 6$~m (SNR of order $10^{2}$--$10^{3}$ at internal-wave scales), rather than the
$\approx 50$~m that a flat floor misleadingly implies. Rotary (clockwise/counter-clockwise)
wavenumber spectra of the complex shear, computed against this floor, resolve the upward- and
downward-propagating internal-wave shear field.

\subsection{Software engineering and reproducibility (improvements)}

The port also modernizes the pipeline as software:

\begin{itemize}
\item \textbf{Configuration-driven, symmetric with the CTD.} Deployment settings live in
\code{config\_adcp.json}; CLI flags override individual values; the resolved settings are
stamped into product attributes.
\item \textbf{Unit-tested numerics.} The transforms, geometry, cast detection, spectral fit,
and structure-function fit are covered by a fast, data-free test suite (white noise $\to$ zero
turbulent amplitude; a synthetic $k^{-5/3}$ field $\to$ the correct slope; etc.), run
alongside the CTD tests.
\item \textbf{Validated against the published product} (Section 3.3), with the comparison
scripts retained under \code{ww\_sig1000/validation/}.
\item \textbf{Crash-safe output.} Products are written to a temporary file and atomically
renamed into place, so a failed or interrupted write never truncates an existing product, and
the swap succeeds even while a reader (e.g.\ an open analysis notebook) holds the previous
file.
\end{itemize}

\section{Summary of ADCP port improvements}

\begin{table}[h]
\centering
\begin{tabular}{@{}p{0.16\textwidth} p{0.36\textwidth} p{0.40\textwidth}@{}}
\toprule
\textbf{Area} & \textbf{Original MATLAB toolbox} & \textbf{PyWirewalker port} \\
\midrule
Ingest & \code{.mat} export $+$ \code{sort\_file}/\code{merge\_signature} & Direct \code{.ad2cp} read via DOLfYN; bounded-memory streaming \\
Turbulence (spectral) & \code{WWturb\_upward.m} ($+0.28$~dex bias) & \code{ProcessSingleProfile.m} method; validated to $-0.00$~dex / corr 0.994 \\
Turbulence (2nd estimator) & Present but unused & Structure-function $\eps$ in every product; gap-robust in low scattering \\
Shear & --- & Beam-differenced, motion-immune shear $+$ per-bin Doppler SEM \\
Motion correction & Single AHRS-based model & v1/v2 models; buoyant-ascent $+$ sail correction; selectable/auto attitude \\
Quality control & --- & $\eps$--noise-floor coupling; transfer-function shear noise floor; rotary spectra \\
Reproducibility & Script constants & Config-driven, unit-tested, provenance in attrs, atomic writes \\
\bottomrule
\end{tabular}
\caption{Improvements introduced during the MATLAB-to-Python port of the ADCP toolbox.}
\end{table}

\appendix

\section{Example CTD configuration (\code{config\_ctd.example.json})}

The generic template shipped in the repository; copy it to \code{config\_ctd.json} and edit.
Every quantity discussed in Section~2 is set here; paths may use \code{\textasciitilde} and are
overridable by environment variables.

\begin{lstlisting}[style=cfg]
{
  "rsk_file": "/path/to/deployment.rsk",
  "output_dir": "/path/to/output",
  "basename": "NOPP-Aleutians_RBR-Concerto-213752",
  "mooring": "NOPP-Aleutians",
  "instrument": "RBRconcerto3 S/N 213752",
  "latitude": 49.5,
  "longitude": -159.0,
  "atmospheric_pressure_dbar": 10.1325,
  "sampling_hz": 2.0,
  "thermal_mass": { "alpha": 0.04, "beta_per_s": 0.1, "gamma": 1.0 },
  "grid": {
    "l2_dz_m": 0.5, "zmin_m": 0.0, "zmax_m": 500.0,
    "l3_dz_m": 1.0, "l3_dt": "30min", "l3_interp_max_gap_bins": 1
  },
  "n2_vertical_smoothing_m": 5.0,
  "gravity": 9.81
}
\end{lstlisting}

\section{Example ADCP configuration (\code{config\_adcp.example.json})}

The generic template shipped in the repository; copy it to \code{config\_adcp.json} and edit.
Every quantity discussed in Section~3 is set here; the instrument geometry is not, since it is
read from the \code{.ad2cp}.

\begin{lstlisting}[style=cfg]
{
  "ad2cp_file": "/path/to/deployment.ad2cp",
  "output_dir": "/path/to/output",
  "basename": "TLC_23_sig1000",
  "mooring": "TLC_23",
  "instrument": "Nortek Signature1000 S/N 101913",
  "latitude": null,
  "longitude": null,
  "velocity":   { "boxsize_m": 1.0, "z_max_m": null, "motion_correct": true },
  "turbulence": { "dep_res_m": 3.0, "max_dep_m": 100.0 },
  "cast":       { "kind": null, "min_span_dbar": 40.0, "corr_min": 50, "chunk": 500000 }
}
\end{lstlisting}

\noindent Two groups of keys are optional and take defaults when absent: the velocity
attitude/motion-model selection (Section~3.2) and the record-trim bounds (Section~3.1). They
are added when a deployment needs them --- for example, the NOPP-Aleutians Signature1000 uses
the buoyant-ascent motion model and trims deployment/recovery transit by time:

\begin{lstlisting}[style=cfg]
  "velocity": { "boxsize_m": 1.0, "motion_correct": true, "motion": "v2", "attitude": "auto" },
  "cast":     { "min_span_dbar": 40.0, "corr_min": 50, "chunk": 500000,
                "start_time": "2024-04-16T00:30:00", "end_time": "2024-10-28T22:30:00" }
\end{lstlisting}

\begin{thebibliography}{9}
\bibitem{lueck1990} Lueck, R.~G. and Picklo, J.~J. (1990). Thermal inertia of conductivity
cells: Observations with a Sea-Bird cell. \emph{J. Atmos. Oceanic Technol.}, 7, 756--768.
\bibitem{northcott2026} Northcott, D. et al. (2026). Wirewalker Signature1000 turbulence
dataset. Dryad, \href{https://doi.org/10.5061/dryad.8sf7m0d44}{doi:10.5061/dryad.8sf7m0d44}
(includes \code{ProcessSingleProfile.m}).
\bibitem{teos10} McDougall, T.~J. and Barker, P.~M. (2011). \emph{Getting started with TEOS-10
and the Gibbs Seawater (GSW) Oceanographic Toolbox.}
\bibitem{wiles2006} Wiles, P.~J. et al. (2006). A novel technique for measuring the rate of
turbulent dissipation in the marine environment. \emph{Geophys. Res. Lett.}, 33, L21608.
\end{thebibliography}

\end{document}
